\paragraph{\trFunctor\trCategory。}\label{sub:functor-categories}
设 $\scr{A}$ 是\trStrictlyUnital{}\trACategory，则对任意\trNonunital{}\trACategory{}%
$\scr{C}$，$\mathit{nu}\text{-}\mathit{fun}(\scr{C},\scr{A})$ 仍是\trStrictlyUnital{}的。
显式地说，严格\trIdentityMorphism{}即那些\trNaturalTransformation{} $E_{\scr{F}}$，其
定义为 $E^0_{\scr{F}}=e_{\scr{F}(X)}
\in\hom^0_{\scr{A}}(\scr{F}(X),\scr{F}(X))$\footnote{译注：原书此处印作
$\hom^0_{\scr{B}}(\scr{F}(X),\scr{F}(X))$；疑应为
$\hom^0_{\scr{A}}(\scr{F}(X),\scr{F}(X))$（$\scr{F}(X)$ 是 $\scr{A}$ 的\trObject，且
本段此前并未出现 $\scr{B}$）；此处径改。——译者注}，而对一切 $d>0$ 有
$E^d_{\scr{F}}=0$。稍欠显然的是下面这条：
\vkthmsync{3}
\begin{lemma}\label{thm:functor-category-cunital}
若 $\scr{A}$ 是\trCunital{}的，则 $\scr{Q}=\mathit{nu}\text{-}\mathit{fun}(\scr{C},\scr{A})$
亦然。此外，若 $E_{\scr{F}}\in\hom^0_{\scr{Q}}(\scr{F},\scr{F})$ 是在 $H(\scr{Q})$ 中
代表\trIdentityMorphism{}的任意一条链，那么对每个 $X\in\Ob\scr{C}$，
$E^0_{\scr{F}}\in\hom^0_{\scr{A}}(\scr{F}(X),\scr{F}(X))$\footnote{译注：原书此处印作
$\hom^0_{\scr{Q}}(\scr{F}(X),\scr{F}(X))$；疑应为
$\hom^0_{\scr{A}}(\scr{F}(X),\scr{F}(X))$（$\scr{F}(X)$ 并非 $\scr{Q}$ 的\trObject，
且本句结论即在 $H(\scr{A})$ 中代表恒等）；此处径改。——译者注}就在 $H(\scr{A})$ 中
代表\trIdentityMorphism。
\end{lemma}
\begin{proof}
这可以化归为\trStrictlyUnital{}的情形，如下。取\trFormalDiffeomorphism{} $\Phi$，满足
$\Phi^1=\id$，使得 $\tilde{\scr{A}}=\Phi_{*}\scr{A}$ 是\trStrictlyUnital{}的，并令
$\tilde{\scr{Q}}=\mathit{nu}\text{-}\mathit{fun}(\scr{C},\tilde{\scr{A}})$。考察与
$\Phi$ 作左\trComposition{}的那个\trAFunctor，并记
$\tilde{\scr{F}}=\Phi\circ\scr{F}$。在\trNaturalTransformation{}的层面上，这给出交换图
\begin{equation}\label{eq:functor-category-diagram}
\begin{tikzcd}
\Hom_{H(\scr{Q})}(\scr{F},\scr{F})
  \ar[rr,"H(\scr{L}^1_{\Phi})"]
  \ar[dr]
&& \Hom_{H(\tilde{\scr{Q}})}(\tilde{\scr{F}},\tilde{\scr{F}})
  \ar[dl] \\
& \Hom_{\mathit{Nu}\text{-}\mathit{Fun}(C,A)}(F,F).
\end{tikzcd}
\end{equation}
这里 $C=H(\scr{C})$、$A=H(\scr{A})=H(\tilde{\scr{A}})$、
$F=H(\scr{F})=H(\tilde{\scr{F}})$ 是\trCohomology{}层面上的结构，而两条斜向下的箭头是
映射\eqref{eq:H-of-nu-fun-to-Nu-Fun}，它们把 \trAFunctor{}之间的%
\trNaturalTransformation{}映到其\trCohomology{}层面的对应物。
\cref{thm:left-composition-cff} 表明，$H(\scr{L}^1_{\Phi})$ 是一个\trIsomorphism。选取
某条链 $E_{\scr{F}}$，使得 $[\scr{L}^1_{\Phi}(E_{\scr{F}})]$ 是\trStrictlyUnital{}元
所在的类 $[E_{\tilde{\scr{F}}}]$。由 $H(\scr{L}^1_{\Phi})$ 的\trFunctoriality，这个类
是 $H(\scr{Q})$ 中的一个\trIdentityMorphism。在
\eqref{eq:functor-category-diagram} 中作图追踪可知，
$H(E_{\tilde{\scr{F}}})=H(E_{\scr{F}})$ 是
$\mathit{Nu}\text{-}\mathit{Fun}(C,A)$ 中的恒等\trNaturalTransformation。这就蕴含：
每一个 $[E^0_{\tilde{\scr{F}}}]=[E^0_{\scr{F}}]$ 都是 $A$ 中的\trIdentityMorphism，
正如所要证。
\end{proof}

\begin{lemma}\label{thm:left-composition-cunital}
设 $\scr{G}\colon\scr{A}\to\scr{B}$ 是\trCunital{}的\trAFunctor。那么，对任意%
\trNonunital{}\trACategory{} $\scr{C}$，\trComposition\trFunctor{}%
$\scr{L}_{\scr{G}}\colon\mathit{nu}\text{-}\mathit{fun}(\scr{C},\scr{A})
\to\mathit{nu}\text{-}\mathit{fun}(\scr{C},\scr{B})$ 是\trCunital{}的。（对右%
\trComposition{}也有一个类似的结果：设 $\scr{G}\colon\scr{A}\to\scr{B}$ 是%
\trNonunital{}\trACategory{}之间的一个\trNonunital{}\trAFunctor。那么，对任意%
\trCunital{}\trACategory{} $\scr{C}$，\trComposition\trFunctor{}%
$\scr{R}_{\scr{G}}\colon\mathit{nu}\text{-}\mathit{fun}(\scr{B},\scr{C})
\to\mathit{nu}\text{-}\mathit{fun}(\scr{A},\scr{C})$ 是\trCunital{}的。）
\end{lemma}
\begin{proof}
设 $E_{\scr{F}}\in\hom^0_{\mathit{nu}\text{-}\mathit{fun}(\scr{C},\scr{A})}
(\scr{F},\scr{F})$ 是在 $H(\mathit{nu}\text{-}\mathit{fun}(\scr{C},\scr{A}))$ 中代表%
\trIdentityMorphism{}的一条链。取它在\trComposition{}下的像
\begin{equation}\label{eq:left-composition-unit}
[\scr{L}^1_{\scr{G}}(E_{\scr{F}})]
\in\Hom^0_{H(\mathit{nu}\text{-}\mathit{fun}(\scr{C},\scr{B}))}
(\scr{G}\circ\scr{F},\scr{G}\circ\scr{F}).
\end{equation}
我们已经看到，对每个 $X$，
$E^0_{\scr{F}}\in\hom^0_{\scr{A}}(\scr{F}(X),\scr{F}(X))$ 代表\trCohomology{}层面的
恒等，因此
$\scr{L}^1_{\scr{G}}(E_{\scr{F}})^0=\scr{G}^1(E^0_{\scr{F}})
\in\hom^0_{\scr{B}}(\scr{G}(\scr{F}(X)),\scr{G}(\scr{F}(X)))$
亦然。利用\cref{thm:right-composition-iso} 以及
$\mathit{nu}\text{-}\mathit{fun}(\scr{C},\scr{B})$ 的\trCunital{}性，可知
\eqref{eq:left-composition-unit} 是一个\trInvertible{}\trMorphism。由
$\scr{L}_{\scr{G}}$ 的\trFunctoriality，它还是\trIdempotent{}的，从而它等于恒等
（这一类论证在后文会被反复使用）。对 $\scr{R}_{\scr{G}}$ 亦有同样的做法。
\end{proof}
